\begin{name}
	{\tenchude}
	{TOÁN 10}
	{LỚP TOÁN THẦY PHÁT}
	{Thời gian: 90 phút - Không kể thời gian phát đề}
\end{name}
\Opensolutionfile{ansbook}[ans/ansbookDe1]
\TN
\Opensolutionfile{ans}[ans/ansDe1-TN1]
\begin{ex}%[0D3H1-5]
	Cho hàm số $ f(x) = \sqrt{x^2 - 2x + 5} $. Mệnh đề nào sau đây là mệnh đề đúng?
	\choice
	{Hàm số nghịch biến trên khoảng $(1;3)$}
	{\True Hàm số đồng biến trên khoảng $(2;3)$}
	{Hàm số nghịch biến trên khoảng $(1;2)$}
	{Hàm số nghịch biến trên khoảng $(2;3)$}
	\loigiai{
		Điều kiện xác định của hàm số là $x^2-2x+5\geq 0$ (luôn đúng với mọi $x$).\\
		Tập xác định của hàm số là $\mathscr{D}=\mathbb{R}$.\\
		Với $x_1$, $x_2$ $\in (1;+\infty)$ và $x_1 < x_2$, ta có
		\allowdisplaybreaks
		\begin{eqnarray*}
			f\left(x_1\right)-f\left(x_2\right)&=& \sqrt{x_1^2 - 2x_1 + 5}-\sqrt{x_2^2 - 2x_2 + 5}\\
			&=& \dfrac{x_1^2-2x_1+5-\left(x_2^2-2x_2+5\right)}{\sqrt{x_1^2 - 2x_1 + 5}+\sqrt{x_2^2 - 2x_2 + 5}}\\
			&=& \dfrac{\left(x_1-x_2\right)\left(x_1+x_2-2\right)}{\sqrt{x_1^2 - 2x_1 + 5}+\sqrt{x_2^2 - 2x_2 + 5}}>0
		\end{eqnarray*}
		Vậy hàm số luôn đồng biến trên $(1;+\infty)$ hay hàm số đồng biến trên khoảng $(2;3)$.
	}
\end{ex}

\begin{ex}%[0D3N2-1]
	Cho hàm số $y=x^2-4x+3$ có đồ thị là một parabol $(P)$. Tìm tọa độ đỉnh $S$ của $(P)$.
	\choice
	{$S(-2;1)$}
	{$S(-2;-1)$}
	{\True $S(2;-1)$}
	{$S(2;3)$}
	\loigiai{
		Tọa độ đỉnh của parabol $y=ax^2+bx+c$ là $S\left(-\dfrac{b}{2a};-\dfrac{\Delta}{4a}\right)$, với $\Delta =b^2-4ac$.\\
		Trong trường hợp này, $a=1$, $b=-4$, $c=3$.\\
		Vậy $\heva{&x_S=-\frac{-4}{2(1)}=2\\&y_S=-\frac{(-4)^2-4\cdot 1\cdot 3}{4\cdot 1}=-\frac{16-12}{4}=-1.}$\\
		Vậy tọa độ đỉnh là $S(2;-1)$.
	}
\end{ex}

\begin{ex}%[0-TK-HK1-CT-2-2425]%[VN-MT-9-CT-10-11, Võ Nguyên Thạch]%[0D3H2-2]
	Cho hàm số $y=x^2-2x-1$, mệnh đề nào đúng?
	\choice
	{Hàm số nghịch biến trên khoảng $(1;+\infty)$}
	{Hàm số đồng biến trên khoảng $(-\infty;1)$}
	{\True Đồ thị hàm số có trục đối xứng $x=1$}
	{Đồ thị hàm số nhận $I(-2;1)$ làm đỉnh}
	\loigiai{
		Xét hàm số $y=x^2-2x-1$.
		\begin{itemize}
			\item Bảng biến thiên
			      \begin{center}
				      \begin{tikzpicture}
					      \tkzTabInit[nocadre=true, espcl=4, lgt=1.0, deltacl=0.5]
					      {$x$/0.7,$y$/2}
					      {$-\infty$,$1$,$+\infty$}
					      \tkzTabVar{+/$+\infty$,-/$-2$,+/$+\infty$}
				      \end{tikzpicture}
			      \end{center}
			\item Đồ thị có đỉnh parabol $I(1;-2)$
			\item Hàm số đồng biến trên khoảng $(1;+\infty)$, nghịch biến trên khoảng $(-\infty;1)$.
			\item Đồ thị hàm số có trục đối xứng $x=1$.
		\end{itemize}
	}
\end{ex}

\begin{ex}%[0D6H1-1]
	Kết quả đo chiều dài của một cây cầu được ghi là $152m\pm 0,2m$. Tìm sai số tương đối của phép đo chiều dài cây cầu đó
	\choice
	{\True $\delta_a<0,1316\%$}
	{$\delta_a<1,316\%$}
	{$\delta_a=0,1316\%$}
	{$\delta_a>0,1316\%$}
	\loigiai{
		Sai số tương đối $\delta_a\le \dfrac{0,2}{152}=0,001315789\approx 0,1316\%$.
	}
\end{ex}

\begin{ex}%[VN-MT-9]%[ 0-TK-HK1-KN-2-2425 , Trần Đức Thắng]%[0D6H4-2]
	Mẫu số liệu hình bên cho biết cân nặng của $10$ trẻ sơ sinh (đơn vị: kg).
	\begin{center}
		\begin{tabular}{|c|c|c|c|c|c|c|c|c|c|}
			\hline
			$2{,}977$ & $3{,}155$ & $3{,}920$ & $4{,}236$ & $3{,}387$ & $2{,}593$ & $3{,}270$ & $3{,}412$ & $3{,}813$ & $4{,}042$ \\
			\hline
		\end{tabular}
	\end{center}
	Khoảng biến thiên (đơn vị: kg) của mẫu số liệu trên bằng
	\choice
	{$0{,}384$}
	{$0{,}194$}
	{\True $1{,}643$}
	{$3$}
	\loigiai{
	Trước hết, ta sẽ sắp xếp mẫu số liệu theo thứ tự không giảm
	\begin{center}
		\begin{tabular}{|c|c|c|c|c|c|c|c|c|c|}
			\hline
			$ 2{,}593$ & $ 2{,}977$ & $ 3{,}155 $ & $ 3{,}270 $ & $3{,}387$ & $ 3{,}412 $ & $3{,}813 $ & $ 3{,}920 $ & $4{,}042$ & $ 4{,}236$ \\
			\hline
		\end{tabular}
	\end{center}
	Khoảng biến thiên là $R=4{,}236-2{,}593=1{,}643$ (kg).
	}
\end{ex}

\begin{ex}%[0D0N1-2]
	Gieo hai đồng tiền một lần. Kí hiệu $S$, $N$ lần lượt để chỉ đồng tiền lật sấp, lật ngửa. Xác định biến cố $M$: \lq\lq Hai đồng tiền xuất hiện hai mặt không giống nhau\rq\rq.
	\choice
	{$M = \left\{NN; SS\right\}$}
	{\True $M = \left\{NS; SN\right\}$}
	{$M = \left\{NS; NN\right\}$}
	{$M = \left\{SS; SN\right\}$}
	\loigiai{
		Biến cố $M$: \lq\lq Hai đồng tiền xuất hiện hai mặt không giống nhau\rq\rq\,là $M = \left\{NS; SN\right\}$.
	}
\end{ex}

\begin{ex}%[0D0N1-3]
	Chọn ngẫu nhiên một cuốn truyện từ $27$ cuốn truyện cổ tích khác nhau và $19$ cuốn tiểu thuyết trinh thám khác nhau. Số phần tử của không gian mẫu là
	\choice
	{$50$}
	{$513$}
	{$8$}
	{\True$46$}
	\loigiai{
		Số phần tử của không gian mẫu là $n(\Omega)=27+19=46$.

	}
\end{ex}

\begin{ex}%[0H9N3-1]%[tex hóa đề CK2 - 2025 - Huyên Nguyễn]
	Trong mặt phẳng tọa độ $Oxy$, đường thẳng $\Delta$: $2x-y+2025=0$ có một vectơ pháp tuyến là
	\choice
	{$\overrightarrow{n}=(1;2)$}
	{\True$\overrightarrow{n}=(4;-2)$}
	{$\overrightarrow{n}=(2;1)$}
	{$\overrightarrow{n}=(-2;-1)$}
	\loigiai{Đường thẳng  $\Delta$: $2x-y+2025=0$ có một vectơ pháp tuyến là $\overrightarrow{n}=(4;-2)$.
	}
\end{ex}

\begin{ex}
	Trong mặt phẳng tọa độ $Oxy$, cho đường thẳng $d_1$ có vectơ chỉ phương $\vec{u_1}$ và đường thẳng $d_2$ có vectơ chỉ phương $\vec{u_2}$. Góc giữa hai đường thẳng $d_1$ và $d_2$ được tính bằng công thức
	\choice
	{$\cot (d_1,d_2)=\dfrac{|\vec{u_1}\cdot \vec{u_2}|}{|\vec{u_1}||\vec{u_2}|}$}
	{\True $\cos (d_1,d_2)=\dfrac{|\vec{u_1}\cdot \vec{u_2}|}{|\vec{u_1}||\vec{u_2}|}$}
	{$\sin (d_1,d_2)=\dfrac{|\vec{u_1}\cdot \vec{u_2}|}{|\vec{u_1}||\vec{u_2}|}$}
	{$\tan (d_1,d_2)=\dfrac{|\vec{u_1}\cdot \vec{u_2}|}{|\vec{u_1}||\vec{u_2}|}$}
	\loigiai{
		Gọi $\alpha$ là góc giữa hai đường thẳng $d_1$ và $d_2$.\\
		Ta có $\cos \alpha = \dfrac{|\vec{u_1}\cdot \vec{u_2}|}{|\vec{u_1}||\vec{u_2}|}$.
	}
\end{ex}

\begin{ex}%[10-11EX-GK2 24-25]%[Lương Như Quỳnh]%[0H9N3-3]
	Trong mặt phẳng $Oxy$, đường thẳng $d\colon x-2y-1=0$ song song với đường thẳng có phương trình nào sau đây?
	\choice
	{$-x+2y+1=0$}
	{\True $-2x+4y-1=0$}
	{$x+2y+1=0$}
	{$2x-y=0$}
	\loigiai{
		Xét đường thẳng $d\colon x-2y-1=0$ và đường thẳng $\Delta \colon -2x+4y-1=0$.\\
		Ta có $\dfrac{1}{-2}=\dfrac{-2}{4}\neq \dfrac{-1}{-1} $.\\
		Suy ra $d\parallel \Delta$.
	}
\end{ex}

\begin{ex}%[0H9N5-1]
	Một tiêu điểm của elip $(E)\colon \dfrac{x^2}{5}+\dfrac{y^2}{4}=1$ là
	\choice
	{$F_2=(0;1)$}
	{\True $F_2=(1;0)$}
	{$F_1=(-3;0)$}
	{$F_1=(-1;-2)$}
	\loigiai{
		Ta có $a^2=5$, $b^2=4 \Rightarrow c^2=a^2-b^2=1 \Rightarrow c=1 \Rightarrow F_1=(-1;0), F_2=(1;0)$.
	}
\end{ex}

\begin{ex}%[0H9N5-5]
	Phương trình nào sau đây là phương trình chính tắc của hypebol?
	\choice
	{$x^2+\dfrac{y^2}{3^2}=1$}
	{$\dfrac{x^2}{16}-y^2=-1$}
	{$\dfrac{x^2}{25}-\dfrac{y^2}{9}=-1$}
	{\True $x^2-\dfrac{y^2}{2}=1$}
	\loigiai{
		Phương trình
		$x^2-\dfrac{y^2}{2}=1 \Leftrightarrow \dfrac{x^2}{1}-\dfrac{y^2}{2}=1$ là phương trình chính tắc của hypebol.
	}
\end{ex}
\Closesolutionfile{ans}

\TNTF
\Opensolutionfile{ans}[ans/ansDe1-TN2]
\begin{ex}%[Dự án EX-10-11-Chuẩn hóa]%[Hoàng Thanh Phương]%[0D6H4-3]
	Mẫu số liệu sau cho biết số ghế trống của một rạp chiếu phim trong 11 ngày
	\begin{center}
	\begin{tabular}{|l|l|l|l|l|l|l|l|l|l|l|}
	\hline 0 & 7 & 8 & 22 & 20 & 15 & 18 & 19 & 13 & 11 & 39 \\
	\hline
	\end{tabular}
	\end{center}
	Các khẳng định sau là đúng hay sai?
	\choiceTF
	{\True Khoảng tứ phân vị của mẫu số liệu là $12$}
	{Không có giá trị bất thường trong mẫu số liệu nêu trên}
	{Quy tròn số trung bình của mẫu số liệu với độ chính xác $d=0.5$ là $a=15$}
	{\True Sai số tương đối của số quy tròn $a$ xấp xỉ $2,27\%$}
	\loigiai{
	Sắp xếp mẫu số liệu theo thứ tự không giảm ta được
	\[0\ 7\ 8\ 11\ 13\ 15\ 18\ 19\ 20\ 22\ 39\]
	Ta thấy $Q_1=8, Q_2=15, Q_3=20$.\\
	\begin{itemchoice}
	\itemch Khoảng tứ phân vị $\Delta_Q=20-8=12$.
	\itemch Ta có $Q_1-\dfrac{3}{2}\Delta_Q=-10$, $Q_3+\dfrac{3}{2}\Delta_Q=38$.\\
	Vậy $39$ là giá trị bất thường của mẫu số liệu trên.
	\itemch Số trung bình cộng của mẫu số liệu là
	\[\overline{x}=\dfrac{0+7+8+11+13+15+18+19+20+22+39}{11}=\dfrac{172}{11}\approx 15,6363.\]
	Vậy quy tròn số trung bình cộng với độ chính xác $d=0,5$ là $16$.
	\itemch Sai số tương đối của số quy tròn $a$ là
	\[\delta_a=\dfrac{|a-\overline{x}|}{a}=\dfrac{|16-\dfrac{172}{11}|}{16} \approx 2,27\%.\]
	\end{itemchoice}
	}
	\end{ex}

\begin{ex}%[0H9H3-2]
	Cho hai điểm $A(3;-3)$, $B(-1;-5)$ và đường thẳng $(d)\colon 4x-3y-2=0$.
	\choiceTF
	{\True Một vectơ chỉ phương của đường thẳng $d$ là $\overrightarrow{u_d}=(3;4)$}
	{Đường thẳng đi qua điểm $A$ và vuông góc với $(d)$ có phương trình $4x+3y=3$}
	{Đường tròn $(C)$ đường kính $AB$ có phương trình là $(x-1)^2+(y+4)^2=10$}
	{\True Đường tròn $(C)$ cắt đường thẳng $(d)$ tại hai điểm phân biệt}
	\loigiai{
		\begin{itemchoice}
			\itemch  \textbf{Đúng}.\\  Vectơ pháp tuyến của  đường thẳng $d$ là $\overrightarrow{n_d}=(4;-3)$ suy ra vectơ chỉ phương  $\overrightarrow{u_d}=(3;4)$.
			\itemch \textbf{Sai}.\\  Đường thẳng vuông góc với đường thẳng $(d)$ nên có dạng: $\Delta\colon 3x+4y+m=0$.\\
			$\Delta$ đi qua $A$ nên thay toạ độ $A$ vào phương trình đường thẳng $\Delta$ ta có \[3\cdot3+4\cdot(-3)+m=0\Leftrightarrow m=3.\]
			Vậy phương trình đường thẳng cần tìm là $\Delta\colon 3x+4y+3=0$.
			\itemch \textbf{Sai}.\\
			Phương trình đường tròn $(C)$ có tâm $I(1;-4)$ và bán kính $R=\dfrac{AB}{2}=\sqrt{5}$.\\
			Vậy phương trình đường tròn $(C)$ là $(x-1)^2+(y+4)^2=5$
			\itemch \textbf{Đúng}.\\
			Ta có $\mathrm{d}(I,(d))=\dfrac{|4\cdot1-3\cdot(-4)-2|}{\sqrt{4^2+(-3)^2}}=\dfrac{14}{5}>\sqrt{5}$.\\
			Vậy đường tròn $(C)$ cắt đường thẳng $(d)$ tại hai điểm phân biệt.
		\end{itemchoice}
	}

\end{ex}
\Closesolutionfile{ans}

\TNSA
\Opensolutionfile{ans}[ans/ansDe1-TN3]
\begin{ex}%[0D3H2-3]
	Đồ thị hàm số $y=5x^2-2x-8$ cắt tia $Ox$ tại điểm có hoành độ bằng bao nhiêu (làm tròn đến hàng phần trăm)?
	\shortans{$1{,}48$}
	\loigiai{
		Ta có $5x^2-2x-8=0\Leftrightarrow \hoac{&x=\dfrac{1+\sqrt{41}}{5} \\&x=\dfrac{1-\sqrt{41}}{5}}$.\\
		Vì giao điểm thuộc tia $Ox$ nên hoành độ giao điểm là dương, do đó hoành độ giao điểm là $x=\dfrac{1+\sqrt{41}}{5}\approx 1{,}48$.
	}
\end{ex}

\begin{ex} %[0D7V2-7]
	Một quả bóng được bắn thẳng lên từ độ cao $2\mathrm{m}$ với vận tốc đầu là $30\mathrm{m} / \mathrm{s}$. Khoảng cách của bóng so với mặt đất sau $t$ giây cho bởi hàm số $h(t)=-4{,}9 t^2+30 t+2$, với $h(t)$ tính bằng đơn vị mét . Quả bóng nằm ở độ cao trên $40\mathrm{m}$ trong thời gian bao lâu? Làm tròn kết quả đến hàng phần trăm.
	\shortans{$2{,}54$}
	\loigiai{
	Quả bóng nằm ở độ cao trên $40\mathrm{m}$ nghĩa là $h(t)>40$.\\
	Điều này tương đương với $-4{,}9 t^2+30 t+2>40\Leftrightarrow -4{,}9 t^2+30 t-38>0$.
	Tam thức bậc hai $f(t)=-4,9 t^2+30 t-38$ có $a=-4{,}9<0$ và $\Delta'=15^2-(-4{,}9) \cdot(-38)=\dfrac{194}{5}>0$.\\
	Do đó $f(t)$ có hai nghiệm phân biệt $t_1 \approx 4{,}33$ và $t_2 \approx 1{,}79$. Khi đó ta có bảng xét dấu
	\begin{center}
		\begin{tikzpicture}
			\tkzTabInit[deltacl=0.5,espcl=2.5,lgt=2,nocadre]
			{$t$/1,$f(t)$/1}
			{$-\infty$,$1{,}79$,$4{,}33$,$+\infty$}
			\tkzTabLine{,-,0,+,0,-,}
		\end{tikzpicture}
	\end{center}
	Suy ra $f(t)$ dương khi $t$ thuộc khoảng $(1{,}79 ; 4{,}33)$.\\
	Vậy quả bóng nằm ở độ cao trên $40\mathrm{m}$ trong $4{,}33-1{,}79=2{,}54\mathrm{s}$.
	}
\end{ex}

\begin{ex}%[0-HK1-CT-4-THTHDHSG-HCM-2324]%[VN-MT-6, VM026]%[0D6H3-2]
	Thống kê độ tuổi kiếm được việc làm đúng chuyên ngành sau khi tốt nghiệp đại học của một nhóm người, thu được bảng sau
	\begin{center}
		\begin{tabular}{|c|c|c|c|c|c|}
			\hline
			Tuổi     & 22 & 23 & 24 & 25 & 26 \\
			\hline
			Số người & 12 & 20 & 19 & 8  & 5  \\
			\hline
		\end{tabular}
	\end{center}
	Tính độ tuổi trung bình của nhóm người này khi họ kiếm được việc làm đúng chuyên ngành (kết quả làm tròn đến chữ số hàng phần mười).
	\shortans[]{$23{,}6$}
	\loigiai{
		Ta có, cỡ mẫu $n=12+20+19+8+5=64$.\\
		Độ tuổi trung bình của nhóm người này khi họ kiếm được việc làm đúng chuyên ngành được tính là
		\[\overline{x}=\dfrac{12\cdot 22 + 20 \cdot 23 + 19 \cdot 24 + 8 \cdot 25+5 \cdot 26}{64}=\dfrac{1510}{64}\approx23{,}6.\]
	}
\end{ex}

\begin{ex}%[0H9V4-2]
	Trong mặt phẳng $Oxy$, gọi $(C)$ là đường tròn đi qua hai điểm $A( 1;1)$, $B(5;3)$ và có tâm $I$ thuộc trục hoành. Hoành độ của điểm $I$ bằng bao nhiêu? \shortans[4]{$4$}
	\loigiai{ Tâm $I$ của đường tròn thuộc trục hoành nên $I(a;0)$.\\
		Đường tròn tâm $I(a;0)$ đi qua hai điểm $A( 1;1)$, $B(5;3)$ nên
		\begin{eqnarray*}
			& & IA=IB\\
			&\Leftrightarrow & \sqrt{(1-a)^2+1^2}=\sqrt{(5-a)^2+3^2}\\
			&\Leftrightarrow & 1-2a+a^2+1=25-10a+a^2+9\\
			&\Leftrightarrow & 8a =32\\
			&\Leftrightarrow & a=4.
		\end{eqnarray*}
		Do đó tâm của đường tròn là $I(4;0)$.
	}
\end{ex}

\Closesolutionfile{ans}

\TL
\begin{ex}%[0D3H1-2]
	Tìm tập xác định của hàm số $y=\dfrac{5}{x-1}+\sqrt{3-x}$
	\loigiai{
		Hàm số xác định khi và chỉ khi $\heva{&x-1\ne 0 \\ &3-x \geq 0} \Leftrightarrow \heva{&x\ne 1\\&x \leq 3.}$\\
		Vậy tập xác định của hàm số $\mathscr{D}=\left(-\infty;3\right]\setminus\{1\}$.
	}
\end{ex}


\begin{ex}%[Huỳnh Xuân Tín, dự án giảng 10-11, CTGDPT2018]%[0D7C3-6]
	\immini{
		Ba thành phố $A$, $B$, $C$ nằm tại ba vị trí lập thành một tam giác vuông tại $B$. Khoảng cách $AB$ là $6$ km, khoảng cách $BC$ là $15$ km. Để nối gần hai thành phố $C$ và $A$, trên đường đi từ $C$ đến $B$ người ta xây xựng một bến tàu tại $M$ cho mọi người có thể đi tàu theo đường sông tới $A$. Bác An di chuyển theo tuyến đường mới, đi bộ từ $C$ đến $M$ với vận tốc $10$ km/h và đi tàu từ $M$ tới $A$ với vận tốc $30$ km/h thì hết $1$ giờ $16$ phút tới được $A$. Tính tổng quãng đường bác An đi từ $A$ tới $C$ theo đơn vị ki-lô-mét.
	}{\begin{tikzpicture}[line join=round, line cap = round, >=stealth, scale=.8,font=\footnotesize,transform shape]
			\foreach \x/\y/\z/\g in {0/0/B/-135, 3/0/M/-90, 5/0/C/-45,0/2/A/90}
			\draw[fill=black] (\x,\y) circle(1pt) coordinate (\z) ($(\z)+(\g:3.5mm)$) node{$\z$};
			\draw (M)--(A)--(B)--(C);
		\end{tikzpicture}	}
	% \shortans{$17$}
	\loigiai{
	Đổi $1$ giờ $6$ phút $=$ $\dfrac{31}{30}$ giờ.\\
	Gọi thời gian đi quãng đường $CM$ của bác An là $x$ (giờ, $0 < x < \dfrac{31}{30}$) thì $CM = 10x$ hay $BM = 15 -10x$.\\
	Thời gian đi quãng đường $MA$ của bác An là $\dfrac{31}{30} - x$ nên $AM = 30\left( \dfrac{31}{30} - x \right) = 31 - 30x$.\\
	Theo định lý Py-ta-go ta có $AM^2 = AB^2 + BM^2$ hay
	\[(31 - 30x)^2 = 6^2 + (15-10x)^2 \Rightarrow 800x^2 - 1560x + 700 = 0 \Rightarrow  \hoac{&x = 0{,}7 & &\text{ (thoả mãn)}\\ &x=1{,}25 & & \text{ (loại)}.}\]
	Do đó tổng quãng đường bác An đi được là $CM + AM = 31 - 20x = 17$ km.
	}
\end{ex}

\begin{ex}%[Đề chuẩn hóa Toán 11]%[Nguyễn Trung Kiên]%[0D0C2-3]
	Xếp ngẫu nhiên $10$ học sinh gồm $2$ học sinh lớp 11A, $3$ học sinh lớp 11B và $5$ học sinh lớp 11C thành một hàng ngang. Biết rằng xác suất để trong $10$ học sinh trên không có $2$ học sinh cùng lớp đứng cạnh nhau bằng $\dfrac{11}{a}$. Giá trị của $a$ bằng bao nhiêu?
	% \shortans[]{$630$}
	\loigiai{
	Phép thử \lq\lq Xếp ngẫu nhiên $10$ học sinh gồm $2$ học sinh lớp 11A, $3$ học sinh lớp 11B và $5$ học sinh lớp 11C thành một hàng ngang\rq\rq\, có số phần tử của không gian mẫu là $n(\Omega)=10!$.\\
	Gọi $A$ là biến cố \lq\lq Trong $10$ học sinh trên không có $2$ học sinh cùng lớp đứng cạnh nhau\rq\rq.
	\begin{itemize}
	\item Xếp $5$ học sinh lớp 11C thành một hàng ngang có $5!$ cách. Khi đó với mỗi cách xếp $5$ học sinh lớp 11C sẽ có $6$ khoảng trống gồm $4$ khoảng trống giữa các học sinh lớp 11C và $2$ khoảng trống ở hai đầu.
	\item Xếp $3$ học sinh lớp 11B vào các khoảng trống, sau đó xếp $2$ học sinh lớp 11A vào các khoảng trống còn lại. Dễ thấy không thể xếp đồng thời $2$ học sinh lớp 11B vào $2$ vị trí hai đầu vì khi đó chắc chắn sẽ có ít nhất $2$ học sinh lớp 11C đứng cạnh nhau.\\
	Ta xét các trường hợp sau:
	\begin{itemize}
	\item[TH1:] Không có học sinh 11B ngồi ở vị trí hai đầu hàng của nhóm học sinh lớp 11C đã xếp. Khi đó, xếp $3$ học sinh lớp 11B vào $4$ vị trí trống ở giữa có $\mathrm{A}_4^3$ cách. Ứng với mỗi cách xếp đó, chọn lấy $1$ trong $2$ học sinh lớp 11A xếp vào vị trí trống thứ $4$ (để hai học sinh lớp 11C không được ngồi cạnh nhau), sau đó xếp học sinh lớp 11A còn lại vào $1$ trong $8$ vị trí (không cạnh học sinh lớp 11A đã xếp).\\
	Trường hợp này có $\mathrm{A}_4^3 \cdot 2 \cdot 8 = 384$ khả năng.
	\item[TH2:] Có $1$ học sinh 11B ngồi một trong hai đầu hàng của nhóm học sinh lớp 11C. Khi đó còn $4$ học sinh còn lại ta xếp vào $4$ khoảng trống giữa các học sinh lớp 11C.\\
	Trường hợp này có $\mathrm{C}_3^1 \cdot 2\cdot 4! = 144$ khả năng.
	\end{itemize}
	\end{itemize}
	Vậy số các kết quả thuận lợi của biến cố $A$ là $n(A)=5!\cdot (384+144)= 63\,360$.\\
	Xác suất của biến cố $A$ là $\mathrm{P}(A)=\dfrac{n(A)}{n(\Omega)}=\dfrac{63\,360}{10!} = \dfrac{11}{630}$.
	}
	\end{ex}
\Closesolutionfile{ansbook}

